\section{Precedents: how each Drawmaha blow-up has been beaten before}\label{sec:precedents}

\mkey{plo practice}No one has solved Drawmaha, but each of its three blow-ups (Section~\ref{sec:drawmaha}) has been beaten separately, and the precedents set this project's expectations.

\textbf{The hand-space blow-up: pot-limit Omaha practice.} PLO's 270{,}725 starting hands are the nearest commercial-scale precedent, attacked three ways. \textbf{MonkerSolver} (industry standard since $\sim$2017): MCCFR over a bucketed tree --- no preflop bucketing, postflop buckets from strength tiers $\times$ board texture $\times$ blocker classes, granularity a user dial against RAM and days of runtime. Battle-tested, and its artifact is instructive: hands inside one bucket cannot differentiate, so bluffs concentrate into few combos --- the bucketing fingerprint. \textbf{GTO Wizard AI} solves PLO street-by-street with self-play-trained EV networks replacing future streets, rivers solved exactly (self-reported: Nash distance $<$0.1\% pot on rivers, $\sim$0.14\% average EV loss on turns, seconds per spot). \mnote{The river-exact + NN-elsewhere hybrid is copyable at hobby scale: solve the post-draw street exactly; let the network cover the rest.}\textbf{Vision GTO} serves precomputed abstracted solves as a database --- the delivery model this project's live-inference dashboard deliberately rejects, for the same reason the market moved on: custom spots explode any precomputed menu.

\textbf{The draw-round blow-up: the Ostroumov solver.} The only detailed public account of an equilibrium solver for a real draw game~\cite{ostroumov}: limit 2-7 triple draw, built 2013--15 for high-stakes pros and used to prepare the 2015 Kuznetsov--Ivey match. The full game would not fit, so: handcraft preflop and first-draw play, solve every post-first-draw subtree independently (the way hold'em solvers treat flops), iterate the handcrafted trunk against the computed EVs by hand. One full run: $\sim$1{,}000 CPU cores for 3 days. Its two transferable lessons: subgame decomposition around the draw works, and --- the upgrade its author called the biggest --- \hlight{tracking even ONE discarded card changed strategies materially: discard blockers (you threw the card their draw needs) are real EV, so discard identity must survive in the infoset even though only the count is public.} Snowing --- standing pat on nothing as a bluff --- only exists if the representation can express it. The cautionary twin is \textbf{Poker-CNN} (2015~\cite{pokercnn}), which learned 2-7 triple draw draws and bets by self-play CNN with the opponent's draw count as an input feature --- it works as a \emph{bot}, and its own author flags the gap: no internal consistency, no equilibrium claim. Exactly the Deep-CFR-vs-plain-self-play distinction of Section~\ref{sec:rl}, played out in a draw game. And the oldest anchor still pays rent: Zadeh's 1977 analytical treatment of five-card draw~\cite{zadeh} derives pre- and post-draw bluffing frequencies from pot odds --- closed-form sanity checks a toy-drawmaha CFR run should reproduce.

\textbf{The split-pot blow-up: the Ho recipe.} A Harvard master's thesis~\cite{ho2015} computed an $\varepsilon$-equilibrium for the jam/fold endgame of no-limit Omaha Hi-Lo --- four cards, split pot --- with CFR+ on cloud GPUs, and its methodology chapter is effectively this project's split-pot engineering spec: enumerate all starting hands via a combinatorial index; \textbf{bit-pack the high rank and low rank into one integer so a single comparison settles both halves}; pre-tabulate win/split/scoop outcomes over suit-isomorphic board classes (152{,}620 of them); and put every drop of split-pot logic in the terminal payoff function so \hlight{the CFR core runs unmodified --- the split pot never touches the algorithm, only the leaves.} MonkerSolver's Omaha-8 support confirms the same at commercial scale. One Ho-specific caveat: his all-pairs matchup precomputation squares badly from 4-card to 5-card hands --- lazy or sampled evaluation replaces it at Drawmaha size.

\mkey{nearest neighbors}Two adjacent literatures round out the map. Gin rummy (the AAAI/EAAI challenge~\cite{ginrummy}) is the academic home of draw/discard games with \emph{public} discards --- one entry learns a value function for the combinatorial discard choice while solving the small endgame exactly with MCCFR, a plausible architecture split for mini-drawmaha --- though its information structure is easier (discard identities are public; Drawmaha hides them). And the absence finding bears repeating with its consequence: with no Drawmaha bot anywhere, \emph{every} baseline in the evaluation plan is homegrown, so the checkpoint ladder and rule-based agents of Section~\ref{sec:evaluation} are not conveniences --- they are the entire comparative universe until someone else builds one.
